Real robots expose the distance between algorithms and deployment. A policy that looks promising in a benchmark still has to survive calibration error, asynchronous sensors, frame drops, action latency, and joint limits. This gap drew me from mechanical design into computer science and embodied intelligence. I began by asking how mechanisms move; I now want to build learning-based robot systems that can collect reliable data, learn from it, and return safely to the physical world.

My education has followed that transition. At Harbin Institute of Technology, Weihai, I earned a Bachelor of Engineering in Mechanical Design, Manufacturing and Automation, where mechanics, engineering drawing, programming, and robotics projects trained me to reason about physical systems rather than abstract components alone. In 2025, I entered Harbin Institute of Technology, Shenzhen, for a second bachelor's degree in Computer Science and Technology. I see this not as a change away from mechanics, but as an expansion of it: software systems, algorithms, data, and machine learning give me the tools to make machines more adaptive. This background has made me most interested in autonomy problems where model performance depends on the full robot system around it.

The clearest example is my recent real-robot VLA/OpenPI project. I built a ROS2 closed loop around a Dobot CR5 manipulator, an Orbbec Astra2 RGB-D camera, and an electric gripper, covering pose servoing, gripper control, teleoperation recording, eye-to-hand calibration, data collection, model adaptation, and real-machine inference. Instead of treating data collection as a file-format task, I designed an HDF5 synchronization pipeline using RGB image timestamps as the master clock. Robot states, target actions, and gripper feedback were aligned to the image stream through binary search and linear interpolation. When early recordings showed severe frame drops, I optimized the ROS2 communication path with Fast DDS shared memory and RGB-only capture, reducing frame drops from 27.15\% to the 2.57\%--4.74\% range.

The same project taught me that robot learning requires engineering discipline before and after training. I built data quality and safety checks for schema completeness, frame drops, pose jumps, workspace bounds, start-pose validation, and slow replay before converting HDF5 episodes into LeRobot v2.0 datasets. I then adapted OpenPI to the CR5 platform by defining a seven-dimensional state/action mapping, training configuration, WebSocket policy serving, and a CR5 inference client. During deployment, I found that maximizing apparent command frequency was not necessarily safer or better: non-blocking execution could make the policy infer from intermediate robot states before the previous motion was completed, causing unstable trajectories. I therefore adopted action-chunk replanning with blocking short-horizon execution, a conservative design that made latency, state feedback, and safety boundaries easier to reason about. This experience shaped my view of embodied AI as a systems problem as much as a modeling problem.

My undergraduate graduation project gave me an earlier foundation in integrating perception, kinematics, and robot visualization. I developed a desktop HMI for weld-seam recognition and a robot digital twin, combining C++/Qt, PyTorch, OpenGL, and analytical kinematics. On the vision side, I compared a baseline U-Net, VGG-U-Net, and a VGG16-pretrained variant, improving weld-seam segmentation accuracy from 64\% to 96.8\% through transfer learning, data augmentation, and dataset expansion. I integrated image selection, Python inference calls, and segmentation-result display into the Qt HMI. On the robotics side, I modeled the LeArm manipulator with modified DH parameters, derived and verified forward and inverse kinematics, and used OpenGL transforms to render STL link models with slider control and mouse interaction. This project made concrete to me that perception becomes useful only when connected to geometry, interfaces, and operator workflows.

At Xbotics, I moved closer to learning-based robotics environments through an internship on MotrixLab, Isaac Lab, quadruped navigation, and reinforcement learning. I migrated Isaac Lab's ANYmal-C navigation task into a MotrixLab NumPy environment, rebuilding reset and step logic, command sampling, observation assembly, reward calculation, and termination checks while preserving the 12-dimensional joint action space, 54-dimensional policy observation, and target position/yaw interface. For rough-terrain navigation, I implemented MuJoCo heightfield height and slope queries, handled coordinate-frame correction, filtered spawn and goal points by elevation and slope, aligned spawn height, and added boundary termination. I also added stability rewards and regression tests. This work taught me that observation definitions, reward terms, terrain assumptions, and tests can determine whether a learned behavior is meaningful or only a simulation artifact.

I also value the engineering culture behind reliable robot systems. In the HERO RoboMaster Team at HIT Weihai, I contributed to the ROS2 migration of a vision framework, separating image acquisition, visual processing, and communication into node-based modules. I worked on the host-MCU communication link, ROS2 topic exchange, and time synchronization so that vision outputs and control messages followed a consistent timeline. This strengthened my ability to work with interfaces, protocols, and real-time constraints in a collaborative environment.

Through graduate study, I want to develop the theoretical and research depth needed to build more capable embodied intelligence systems. I hope to study robot learning, perception, control, machine learning, distributed robotic systems, and safe real-robot deployment in a more systematic way. My current questions are practical but research-driven: how should real-robot datasets be structured for temporally meaningful actions; how can VLA policies adapt to new manipulators with limited data; and how can deployment pipelines make uncertainty, latency, and safety constraints explicit?
